% TOP_PROBLEM: {"id": 3200008, "title": "Rudin's conjecture on squares in progressions", "category_id": 2, "category": {"id": 2, "name": "combinatorics", "display_name": "Combinatorics", "description": "Counting problems, graph theory, discrete structures.", "slug": "combinatorics", "order_index": 2, "created_at": "2026-07-31T15:26:25.671Z"}, "status": "partially_solved"}
% FIRST_POSTED: null
% ENTRY_KIND: statement_only
% Imported from: batch12.tex; UnsolvedMath ID 3200008
% Compile twice: lualatex 3200008.tex
\documentclass[11pt]{article}
\usepackage{fontspec}
\setmainfont{Latin Modern Roman}
\usepackage{amsmath,amssymb,mathtools,unicode-math}
\setmathfont{Latin Modern Math}
\usepackage{xurl,hyperref}
\hypersetup{hidelinks,pdftitle={UnsolvedMath Problem 3200008}}
\newcommand{\sref}[2]{\hyperlink{source:#1:#2}{[S#2]}}
\newcommand{\cataloguescope}[1]{\paragraph{Mathematical question.}#1}
\begin{document}
% UnsolvedMath ID 3200008; source catalogue code AMR-031-0008
\setcounter{section}{3200007}
\section{Rudin's conjecture on squares in progressions}\label{3200008}
{\small\sffamily UnsolvedMath ID 3200008 · Combinatorics}
\subsection{Definitions and mathematical statement}

\paragraph{Shared notation.}$\mathbb N=\{1,2,\ldots\}$, and $\mathbb Z,\mathbb Q,\mathbb R,\mathbb C$ denote the integers, rationals, reals and complex numbers. Any other conventions are specified locally.


For positive integers $N,q,a$, define
\[
Q(N;q,a)=\#\{j\in\{0,\ldots,N-1\}: a+jq=t^2\text{ for some integer }t\},
\qquad Q(N)=\max_{q,a\ge1}Q(N;q,a).
\]
This maximum exists because the possible counts belong to the finite set $\{0,\ldots,N\}$. The progression with $q=24$ and $a=1$ gives the comparison value $Q(N;24,1)$.
\paragraph{Statement recorded in the supplied source.}
Prove the bound $Q(N)=O(\sqrt N)$: there should exist an absolute constant $C>0$ such that $Q(N)\le C\sqrt N$ for all $N\ge1$. In the stronger form, prove
\[
Q(N)=Q(N;24,1)\qquad\text{for every integer }N>6.
\]
Both parts concern arithmetic progressions with positive starting term and positive common difference, as in this catalogue record. \sref{3200008}{2}
\subsection{Short English statement}
Bound the number of squares in any $N$-term positive arithmetic progression by a constant times $\sqrt N$. The stronger conjecture says that, for $N>6$, the progression starting at $1$ with difference $24$ attains the maximum.
\subsection{Sources}
\begin{description}
\item[\hypertarget{source:3200008:1}{S1}] ULAM.AI and the UnsolvedMath Contributors, UnsolvedMath: A Curated Collection of Open Mathematics Problems, record 3200008 (AMR-031-0008). CC BY 4.0. \url{https://huggingface.co/datasets/ulamai/UnsolvedMath}.
\item[\hypertarget{source:3200008:2}{S2}] Enrique González-Jiménez and Xavier Xarles, On a conjecture of Rudin on squares in arithmetic progressions (2013), Introduction, pp.1–3. \url{https://arxiv.org/abs/1301.5122}.
\end{description}
\label{end:3200008}
\end{document}
